
\sgaStableTarget{sga6.e10.eq.3-3-1}@@M0@@

\sgaStableTarget{sga6.e10.sec.3.4}\subsection*{3.4.}
Finally, suppose that the morphisms (3.2.3) and (3.3.1) are both defined for a morphism @@M1@@ such that @@M2@@. This is the case, for example, when @@M3@@ is quasi-compact with an ample @@M4@@-module and @@M5@@, where @@M6@@ is a locally free @@M7@@-module of rank @@M8@@. Then, for every @@M9@@ and @@M10@@, the projection formula in the derived category gives the projection formula
@@M11@@
When @@M12@@ is noetherian, 1.3.2 yields the corresponding graded formula.

\sgaStableTarget{sga6.e10.sec.4}\section*{4. Intersection numbers}

\sgaStableTarget{sga6.e10.sec.4.1}\subsection*{4.1.}
Consider a proper scheme @@M13@@ over a field @@M14@@. Composing the product (3.1.1)
@@M15@@
with the Euler--Poincaré characteristic
@@M16@@
gives a pairing called the ``\emph{intersection number}'':
\sgaStableTarget{sga6.e10.eq.4-1-1}@@M17@@

